\section{1882 Poincaré i niesprzeczność} % lang-pl
\label{sekcja_poincare}

Tego, co Beltrami zrobi po cichu, Poincaré dokona głośno i przy okazji zupełnie innej sprawy. % lang-pl
W tej sekcji zobaczymy, skąd Poincaré weźmie swój model, jak on wygląda i co dokładnie z niego wynika: niesprzeczność geometrii Łobaczewskiego względem euklidesowej oraz niezależność piątego postulatu. % lang-pl

Henri Poincaré urodzi się 29 kwietnia 1854 roku w Nancy, a umrze 17 lipca 1912 roku w Paryżu. % lang-pl
Model dysku wykorzysta w 1882 roku w pracach o funkcjach automorficznych (fuchsowskich), w których będzie nieformalnym rywalem Kleina od 1881 roku; szersze uznanie zdobędzie on dopiero po wydaniu w 1905 roku jego filozoficznej książki \emph{La science et l'hypothèse}. % lang-pl
\index[persons]{Poincaré, Henri}%

\begin{definition}
[model Poincarego] % lang-pl
	Ustalmy okrąg $\Sigma$ (\emph{okrąg podstawowy}). % lang-pl
	Punktami są punkty wnętrza $\Sigma$; prostymi -- części wnętrza $\Sigma$ dowolnych „okręgów'' (okręgów lub prostych) prostopadłych do $\Sigma$; długość odcinka $AB$ to $\log(AB, TS)$, gdzie $S$, $T$ są przecięciami „okręgu'' zawierającego $AB$ z $\Sigma$; miara kąta to miara euklidesowa kąta między „okręgami''. % lang-pl
\index{model!Poincarego} % lang-pl
\end{definition}

Eves \cite[s. 347--348]{eves1_1972}. % lang-pl

Dysk ten jest \emph{konforemny} (zachowuje kąty), a ten sam dysk pojawił się już w pracy Beltramiego z 1868 roku. % lang-pl
Kilkadziesiąt lat później M.~C. Escher użyje go w serii \emph{Circle Limit} (1958--1960) po rozmowach z Harold Coxeterem. % lang-pl
\index[persons]{Escher, Maurits Cornelis}%
\index[persons]{Coxeter, Harold}%

\begin{proposition}
	Jeśli geometria euklidesowa (na płaszczyźnie) jest niesprzeczna, to geometria Łobaczewskiego (na płaszczyźnie) też jest niesprzeczna. % lang-pl
\end{proposition}

Eves \cite[s. 347-353]{eves1_1972}. % lang-pl

Eves bierze za punkt wyjścia układ aksjomatów Hilberta, w którym zamienia aksjomat równoległych (IV-1) na aksjomat Łobaczewskiego IV$'$-1: przez punkt spoza prostej przechodzą co najmniej dwie proste, które jej nie przecinają \cite[s. 347]{eves1_1972}. % lang-pl
Sprawdza wszystkie aksjomaty grup I--III (łączności, porządku, przystawania), IV$'$-1 i V$'$-1 (ciągłość) w modelu: większość jest oczywista, a przystawanie wymaga lematu o niezmienniczości długości względem inwersji w okręgu prostopadłym do $\Sigma$ \cite[s. 348--352]{eves1_1972}. % lang-pl

Dowód ma jeszcze drugą stronę: skoro w modelu (będącym częścią geometrii euklidesowej) spełnione są wszystkie postulaty euklidesowe oprócz piątego, a piąty zawodzi, to piąty postulat \emph{nie} wynika z pozostałych -- jest od nich niezależny \cite[s. 352]{eves1_1972}. % lang-pl
Eves zaznacza też, że sama niesprzeczność geometrii euklidesowej sprowadza się ostatecznie do niesprzeczności liczb rzeczywistych, która pozostaje otwartym pytaniem \cite[s. 352--353]{eves1_1972}. % lang-pl

Greenberg \cite[s. 213]{greenberg_2010} podsumowuje: od chwili, gdy pod koniec XIX wieku pokazano modele Beltramiego--Kleina i Poincarego, wiadomo, że jeśli płaska geometria euklidesowa jest niesprzeczna, to hiperboliczna też, a modele te rozwiały wątpliwości co do tej alternatywy. % lang-pl
Wszystkie płaszczyzny hiperboliczne w sensie Hilberta są, z dokładnością do izomorfizmu, modelami Kleina i Poincarego nad ciałami euklidesowymi \cite[s. 209]{greenberg_2010}. % lang-pl
Poincaré napisze po wykazaniu względnej niesprzeczności, że nikt nie wątpi, iż zwykła geometria jest wolna od sprzeczności, ale skąd ta pewność i w jakim stopniu jest uzasadniona, to pytanie bardzo ważne, które wymaga dalszej pracy; Greenberg widzi w nim także pytanie matematyczne, które logicy później rozstrzygną \cite[s. 214]{greenberg_2010}. % lang-pl
Implikację odwrotną pokażą w 1995 roku A.~Ramsey i R.~D. Richtmyer: model płaszczyzny euklidesowej wewnątrz płaszczyzny hiperbolicznej, w którym „prostymi euklidesowymi'' są proste hiperboliczne przez ustalony punkt oraz krzywe równoodległe od nich. % lang-pl
Model ten pokaże, że dawni geometrzy, jak Clavius (zobacz sekcję \ref{sekcja_dobre_stare}), którzy uważali krzywe równoodległe za proste, mieli częściowo rację, o ile pracowali w płaszczyźnie hiperbolicznej \cite[s. 214]{greenberg_2010}. % lang-pl
\index[persons]{Ramsey, Arlan}%
\index[persons]{Richtmyer, Roger D.}%

Model służy też do obliczeń: to, co trudno wyprowadzić w geometrii Łobaczewskiego, często łatwo wyprowadzić z jego odpowiednika w geometrii euklidesowej, a nawet cała trygonometria hiperboliczna wynika z modelu \cite[s. 353 nn.]{eves1_1972}; Eves opublikuje ją razem z V.~E. Hoggattem w \emph{American Mathematical Monthly} w 1951 roku. % lang-pl
\index[persons]{Hoggatt, Verner Emil}%
\index[persons]{Eves, Howard}%


Hartshorne \cite[s. 355--356]{hartshorne_2000} buduje ten sam model nad dowolnym ciałem euklidesowym $F$ i nazywa jego elementy P-punktami i P-prostymi; wszystkie aksjomaty płaszczyzny Hilberta sprawdza po kolei (prop. 39.1--39.9, s. 356--362). % lang-pl
Kluczem jest inwersja względem okręgu: jest to przekształcenie płaszczyzny, którego Grecy nie mieli w zwyczaju rozważać, bo przekształcenia nie zachowujące odległości ani proporcji były im obce \cite[s. 334]{hartshorne_2000}; inwersja w P-prostej daje w modelu ruch sztywny, czyli odbicie. % lang-pl
Definicja przystawania odcinków przez dwustosunek jest tu wprowadzona bez logarytmu: odcinki $AB$ i $A'B'$ są przystające, gdy $(AB, PQ) = (A'B', P'Q')$; funkcja $\mu(AB) = (AB, PQ)^{-1}$ jest multiplikatywną funkcją odległości \cite[s. 358, 363]{hartshorne_2000} (lemat 39.10). % lang-pl
Model spełnia aksjomat Archimedesa dokładnie wtedy, gdy spełnia go ciało $F$ (prop. 39.11), a dla każdego punktu i półprostej istnieje w nim półprosta równoległa graniczna (prop. 39.12) -- czyli spełnia aksjomat hiperboliczny Hilberta \cite[s. 363--364]{hartshorne_2000}. % lang-pl

\begin{theorem}
[wzór Bolyaia] % lang-pl
\index{wzór!Bolyaia (na kąt równoległości)} % lang-pl
	W modelu Poincarego, jeśli $PQ$ jest prostopadłą z $P$ na prostą $l$, a $m$ jest prostą równoległą graniczną do $l$ przez $P$, tworzącą z $PQ$ kąt $\alpha$, to $\tan\tfrac{\alpha}{2} = \mu(PQ)^{-1}$. % lang-pl
\end{theorem}

Dowód (krótki rachunek na współrzędnych): Hartshorne \cite[s. 364--365]{hartshorne_2000} (prop. 39.13). % lang-pl
Wynika z tego, że dla każdego kąta ostrego istnieje odcinek o takim kącie równoległości -- w modelu istnieje więc absolutna jednostka długości, w przeciwieństwie do geometrii euklidesowej, gdzie wybór jednostki jest dowolny \cite[s. 365--366]{hartshorne_2000} (zobacz sekcję \ref{sekcja_trojkaty_graniczne}). % lang-pl

Zadania: suma kątów w każdym trójkącie modelu jest mniejsza od $\pi$, więc model jest półhiperboliczny (zad. 39.2); istnieje w nim trójkąt o dowolnych kątach o sumie mniejszej niż $\pi$ (zad. 39.6) -- Hartshorne \cite[s. 366]{hartshorne_2000}. % lang-pl


